\subsection{Lựa chọn nền tảng và công nghệ}
\label{sec:platform-selection}

\subsubsection{Phương án đã cân nhắc}
\label{sec:platform-comparison}

\begingroup
\small
\renewcommand{\arraystretch}{1.2}
\setlength{\tabcolsep}{3pt}
\begin{longtable}{>{\RaggedRight\arraybackslash}p{0.17\linewidth}>{\RaggedRight\arraybackslash}p{0.28\linewidth}>{\RaggedRight\arraybackslash}p{0.22\linewidth}>{\RaggedRight\arraybackslash}p{0.23\linewidth}}
\caption{Đối chiếu các nền tảng và framework}
\label{tab:platform-comparison}\\
\toprule
\textbf{Phương án} & \textbf{Điểm phù hợp} & \textbf{Chi phí trong bài} & \textbf{Kết luận} \\
\midrule
\endfirsthead
\toprule
\textbf{Phương án} & \textbf{Điểm phù hợp} & \textbf{Chi phí trong bài} & \textbf{Kết luận} \\
\midrule
\endhead
\bottomrule
\endfoot
Android Native/Kotlin & Trực tiếp dùng lifecycle, Media3, media session và audio focus; nhóm có kinh nghiệm & Chỉ triển khai Android & Chọn \\
iOS Native/Swift & Tích hợp hệ thống iOS trực tiếp & Thêm nền tảng và công cụ; không có bản iOS để kiểm thử & Phân tích, không triển khai \\
Flutter & Chia sẻ UI và logic đa nền tảng & Cần kiểm tra plugin và phần tích hợp Android cho từng yêu cầu media & Không chọn trong phạm vi hiện tại \\
React Native & Chia sẻ UI/logic bằng JS/TS & Phụ thuộc thư viện/native integration cho media và background & Không chọn \\
Kotlin Multiplatform & Có thể chia sẻ logic Kotlin & Vẫn cần chiến lược UI/media theo nền tảng; tăng phạm vi kiểm thử & Cân nhắc sau nếu có yêu cầu iOS \\
\end{longtable}
\endgroup

Không dùng bảng này để khẳng định một framework luôn nhanh hơn framework khác hoặc mọi cross-platform app đều cần tự viết native bridge.

\subsubsection{Technology stack và mức triển khai}
\label{sec:tech-stack}

Bảng~\ref{tab:tech-stack} liệt kê công nghệ theo vai trò trong kiến trúc và mức triển khai tại bản nguyên mẫu hiện tại.

\begingroup
\small
\renewcommand{\arraystretch}{1.2}
\setlength{\tabcolsep}{3pt}
\begin{longtable}{>{\RaggedRight\arraybackslash}p{0.23\linewidth}>{\RaggedRight\arraybackslash}p{0.32\linewidth}>{\RaggedRight\arraybackslash}p{0.37\linewidth}}
\caption{Technology stack và mức triển khai}
\label{tab:tech-stack}\\
\toprule
\textbf{Vai trò} & \textbf{Công nghệ chốt} & \textbf{Hiện trạng} \\
\midrule
\endfirsthead
\toprule
\textbf{Vai trò} & \textbf{Công nghệ chốt} & \textbf{Hiện trạng} \\
\midrule
\endhead
\bottomrule
\endfoot
UI/design system & Compose, Material 3, font đóng gói & Có \\
Navigation & Navigation Compose typed routes + Kotlin Serialization & Có \\
State/DI & ViewModel, StateFlow, lifecycle-aware collection, Hilt/KSP & Có \\
Domain & Kotlin/JVM, coroutines/Flow cho contracts & Có \\
Media & Media3 ExoPlayer, adapter thực thi MediaSessionPort & Có phát sample stream thật \\
Structured local data & Room & Có 6 bảng: video progress, watchlist, điểm, nhiệm vụ hằng ngày, âm thanh trích, nội dung tải (schema v5) \\
Preferences & DataStore & Có \\
Remote API & Retrofit + OkHttp + Kotlin Serialization & Chọn cho giai đoạn remote; chưa có client/backend thật \\
Artwork/ảnh & Coil khi tích hợp ảnh từ nguồn thật & Chưa tích hợp; hiện dùng minh họa mock \\
Music background & MediaSessionService + MediaController & Bắt buộc cho hành vi background đích; chưa triển khai \\
Deferred sync & WorkManager & Dành cho sync/retry trì hoãn; chưa có worker \\
Mini game & Logic Kotlin + Compose/Canvas phù hợp game nhỏ & Có: board game và mini-game nhịp chạy thật, tích hợp thưởng điểm \\
Verification & JUnit, Compose UI tests, verifyArchitecture & Có baseline; mỗi yêu cầu mới cần bằng chứng riêng \\
\end{longtable}
\endgroup

Không thêm dependency, empty service hoặc empty use case chỉ để bảng stack trông đầy đủ. Những mục chưa có phải giữ trạng thái dự kiến cho đến khi có code và test.
